% Generated draft: requires mathematical and source review.
\documentclass[11pt]{article}
\usepackage[margin=25mm]{geometry}
\usepackage{fontspec}
\setmainfont{Latin Modern Roman}
\usepackage{amsmath,amssymb,mathtools,unicode-math}
\setmathfont{Latin Modern Math}
\usepackage{enumitem,needspace,xurl,hyperref}
\setlength{\parindent}{0pt}
\setlength{\parskip}{5pt}
\setlength{\emergencystretch}{3em}
\newcommand{\N}{\mathbb{N}}
\newcommand{\Z}{\mathbb{Z}}
\newcommand{\Q}{\mathbb{Q}}
\newcommand{\R}{\mathbb{R}}
\newcommand{\C}{\mathbb{C}}
\newcommand{\F}{\mathbb{F}}
\newcommand{\A}{\mathbb{A}}
\newcommand{\PP}{\mathbb{P}}
\newcommand{\E}{\mathbb{E}}
\newcommand{\eps}{\varepsilon}
\newcommand{\dd}{\,\mathrm{d}}
\begin{document}
\section*{UnsolvedMath ID: 11000155}
\textbf{Source problem code:} catalogue entry
\subsection{Definitions and mathematical statement}
Let \(\operatorname{Map}_{g,n}\) and its positive submonoid \(\operatorname{Map}_{g,n}^+\) have the source's surface and curve conventions. Typically the latter is generated by right-handed Dehn twists in admissible simple closed curves, but the excerpt does not define whether \(n\) counts boundary components or marked points, or which curves are allowed. Fix \(T\in\operatorname{Map}_{g,n}^+\), and let
\[\mathcal F(T)=\coprod_{r\geq0}\{(t_{c_1},\ldots,t_{c_r}):t_{c_1}\cdots t_{c_r}=T,\ \text{all twists positive and admissible}\}.\]
The \(r=0\) convention matters only if the identity is included; the source does not state it. No fixed factorization length is imposed.

A Hurwitz move replaces adjacent factors \((a,b)\) by \((b,b^{-1}ab)\), or its inverse, preserving their product. Conjugation preserves positivity of Dehn twists. Global conjugation sends every factor \(a_i\) to \(h^{-1}a_i h\), changing the product to \(h^{-1}Th\). Thus on the fixed-product set \(\mathcal F(T)\) it must use \(h\in C_{\operatorname{Map}_{g,n}}(T)\), or else the classification must explicitly include all products conjugate to \(T\). The source does not state this convention.

The broad task is to identify those \(T\) for which positive factorizations can be classified. Its concrete subquestions ask for which \(T\) the orbit set under the specified Hurwitz and global-conjugation equivalence has cardinality one, or is finite:
\[|\mathcal F(T)/\!\sim|=1,\qquad\text{or}\qquad|\mathcal F(T)/\!\sim|<\infty.\]
Since Hurwitz moves and conjugation preserve length, these are all-length assertions; no insertion/deletion, stabilization or other factorization relation is added. The source does not prescribe a classification algorithm, invariant list, genus range or particular family of \(T\).
\subsection{Short English statement}
Identify mapping classes whose positive Dehn-twist factorizations admit a classification, especially those with one or finitely many Hurwitz/global-conjugacy classes. The surface, admissible curves and fixed-product conjugation convention must be specified.
\subsection{Sources}
\begin{itemize}
\item Local UnsolvedMath registry, permanent ID \(11000155\), exact source code \texttt{catalogue entry}.
\item Farb, \emph{Problems on Mapping Class Groups and Related Topics}, Question 2.4, PDF page 144: \url{https://www.math.uchicago.edu/~farb/papers/mcgbook.pdf}. Registry access 2026-07-29; pdf-font-structure extraction.
\item The summary mentions Wajnryb, Stipsicz and Korkmaz as related factorization literature, without precise citations or a complete classification theorem.
\end{itemize}
\end{document}
